\section{Finite one-point laws and physical localization}
\label{ll:section}

This section develops the one-point estimates used by both Hausdorff
bounds.  Finite Green changes of law give approach probabilities and
show that the centered radius stops have a finite Loewner-time limit as
their radii tend to zero; the latter fact supplies finite-time mass in
Section~\ref{retain:section}.  A conformal localization estimate then
controls the image of an initial trace uniformly over maps whose domains
may have narrow passages.  It holds for every $0<\kappa<8$ and is an
input to the ribbon construction.  When $4<\kappa<8$, combining it with
a swallowing estimate gives the physical-disk rule used in the mesh tail.
All constants below depend only on the fixed parameter $\kappa$, unless
another dependence is displayed.

Write
\[
 p=1-\frac\kappa8,\qquad a_0=\frac8\kappa-1,
 \qquad 0<\kappa<8.
\]
We use the original capacity normalization
$\partial_tg_t=2/(g_t-\sqrt\kappa B_t)$.
The conformal radius $R_D(z)$ is $|f'(0)|$ for a conformal map
$f:\mathbb D\to D$ with $f(0)=z$.
In particular, Koebe's Theorem gives
\begin{equation}
 \operatorname{dist}(z,\partial D)\le R_D(z)
 \le4\operatorname{dist}(z,\partial D).
 \label{ll:koebe}
\end{equation}
These inequalities require no boundary regularity
\cite[Section~3.2]{beliaev2015}.

\subsection{Stopped Green tilts}

The following centered diffusion is the usual finite-radius Green
change of law; see \cite[Section~2]{lawlerzhou2013} and
\cite[Sections~2.3--2.4]{lawlerwerness2013}.
We include the exhaustion argument because our later stopping times
may follow arbitrarily irregular pasts.

Let $J:\mathbb H\to D$ be conformal.  The boundary parameters $0$ and
infinity specify the starting and target prime ends of $D$; no continuous
boundary extension of $J$ is assumed here.  Let $K_t$ and $g_t$ be the
filled hulls and map-out functions of ordinary half-plane SLE, and set
$D_t=J(\mathbb H\setminus K_t)$.  Throughout this subsection, $t$ is
Loewner time in the initial half-plane, where the hull has capacity $2t$;
it need not be a capacity parametrization intrinsic to $D$.

For $z\in D$, put $\zeta=J^{-1}(z)$ and
$T_z=\inf\{t:z\notin D_t\}$, with $T_z=\infty$ if this set is empty.
We call the point live at times $t<T_z$.  Until $T_z$, set
\[
 Z_t=g_t(\zeta)-\sqrt\kappa B_t=X_t+iY_t,
 \quad\theta_t=\arg Z_t,
 \quad R_t=R_{D_t}(z)=\frac{2Y_t|J'(\zeta)|}{|g_t'(\zeta)|}.
\]
Define the radius clock by $s=\log(R_0/R_t)$.
Differentiation and It\^o's Formula give
\begin{align}
 ds&=\frac{4Y_t^2}{|Z_t|^4}\,dt,
 &\mathcal L&=\frac\kappa8\partial_\theta^2
       +\frac{\kappa-4}{4}\cot\theta\,\partial_\theta,
 \label{ll:clock}\\
 M_t(z)&=R_t^{-p}\sin^{a_0}\theta_t,
 &dM_t&=\sqrt\kappa a_0\frac{X_t}{|Z_t|^2}M_t\,dB_t.
 \label{ll:density}
\end{align}
The sign of the last Brownian motion depends only on the driver
convention and has no later effect. For the displayed driver it is
the positive sign written here.
In particular,
\begin{equation}
 |\text{Brownian coefficient of }M_t|
 \le C R_t^{-p}\frac{\sin^{a_0}\theta_t}{|Z_t|}.
 \label{ll:coefficient}
\end{equation}
The density is continued by zero after killing.

\begin{lemma}[Finite stopped centered law]\label{ll:finite-tilt}
For $0\le u<\infty$, let $\tau_u$ be the first time the live point has
$R_t=R_0e^{-u}$, with value infinity if this level is never reached.
Then
\[
 \mathbb E\left[M_{\tau_u}(z)
                   \mathbf1_{\{\tau_u<\infty\}}\right]=M_0(z).
\]
The normalized weight defines a consistent law on the stopped path
through $\tau_u$. Under that law, in radius time,
\begin{equation}
 d\theta_s=\frac{\sqrt\kappa}{2}\,dW_s
                    +\cot\theta_s\,ds,
 \qquad
 \mathcal L^*=\frac\kappa8\partial_\theta^2
                    +\cot\theta\,\partial_\theta.
 \label{ll:tilt-generator}
\end{equation}
Both endpoints are unattainable at finite radius time. These claims
hold conditionally after any finite stopping time at which the point
is live, with the same constants and the restarted conformal map.
\end{lemma}

\begin{proof}
First stop at $\tau_u$ or when $\theta$ leaves $(1/N,\pi-1/N)$,
$Y$ falls below $1/N$, or Loewner time reaches $N$.  The stochastic
logarithm coefficient is bounded there, since
\[
 \frac{|X_t|}{|Z_t|^2}\le\frac1{2Y_t}\le\frac N2.
\]
Girsanov's Theorem therefore applies. Its added
angular drift is $(\kappa a_0/4)\cot\theta$; together with
\eqref{ll:clock} this gives \eqref{ll:tilt-generator}.

For $0<b<a_0$, let $V_b(\theta)=\sin^{-b}\theta$. Direct
differentiation gives
\begin{equation}
 \mathcal L^*V_b
 =V_b\left\{\frac{\kappa b}{8}
          -\frac{\kappa b(a_0-b)}8\cot^2\theta\right\}
 \le\frac{\kappa b}{8}V_b.
 \label{ll:lyapunov}
\end{equation}
The stopped nonnegative supermartingale
$e^{-\kappa bs/8}V_b(\theta_s)$ bounds the probability of leaving
$(\delta,\pi-\delta)$ before $u$ by
$e^{\kappa bu/8}V_b(\theta_0)\sin^b\delta$.
This tends to zero. On every compact interval on which the angle
stays in $(0,\pi)$, the remaining variables satisfy
\begin{equation}
 \frac{d\log Y_s}{ds}=-\frac1{2\sin^2\theta_s},
 \qquad
 \frac{dt}{ds}=\frac{Y_s^2}{4\sin^4\theta_s}.
 \label{ll:physical-clock}
\end{equation}
They keep $Y$ positive and Loewner time finite there. Thus the
localized changed laws lose no probability before $u$ as the
cutoffs are removed. The localized density identities and monotone
convergence give the asserted expectation one. Repeating the same
argument between two finite levels gives consistency. Composition
of the Loewner maps and Brownian strong Markov give the conditional
statement. No change of law on a completed infinite-future sigma
field is involved.
\end{proof}

\begin{lemma}[Uniform inverse-angle moment]\label{ll:inverse-angle}
For every $u_0>0$ there is $C_{u_0}<\infty$ such that
\[
 \sup_{\theta\in(0,\pi)}\sup_{s\ge u_0}
 \mathbb E^*_{\theta}\sin^{-a_0}\theta_s\le C_{u_0}.
\]
In particular, after one unit of radius time the probability that
$\sin\theta$ exceeds a fixed positive constant is bounded below,
uniformly over the starting angle.
\end{lemma}

\begin{proof}
We give a short spectral proof which also avoids any ambiguity about
comparison inequalities for unbounded test functions. The invariant
probability measure for \eqref{ll:tilt-generator} is
\[
 \pi(d\theta)=c_\kappa\sin^{8/\kappa}\theta\,d\theta.
\]
For $v=\cos\theta$, the generator and invariant measure become
\[
 \mathcal A=\frac\kappa8(1-v^2)\partial_v^2
                   -\left(1+\frac\kappa8\right)v\partial_v,
 \qquad
 \varpi(dv)=\widetilde c_\kappa
                     (1-v^2)^{4/\kappa-1/2}\,dv.
\]
Integration by parts makes $\mathcal A$ symmetric on polynomials.
If $P_m$ are the orthonormal polynomials for $\varpi$, with $P_0=1$,
its action on the leading coefficient and symmetry imply
\[
 \mathcal AP_m=-\lambda_mP_m,
 \qquad \lambda_m=\frac\kappa8m^2+m.
\]
For completeness, choose $m+1$ separated intervals of length
$c/(m+1)$ in $[-1/2,1/2]$. Since the weight is bounded below there,
each contains a point where $|P_m|\le C\sqrt{m+1}$.
Lagrange interpolation at these points gives
\[
 \|P_m\|_{\infty,[-1,1]}
          \le \exp\{C(m+1)\log(m+2)\}.
\]
Consequently
\[
 k_s(v,w)=1+\sum_{m\ge1}e^{-\lambda_ms}P_m(v)P_m(w)
\]
converges absolutely and uniformly for $s\ge u_0$, with a uniformly
bounded sum. It\^o's Formula for each bounded polynomial gives
$\mathbb E_vP_m(v_s)=e^{-\lambda_ms}P_m(v)$.
Polynomial density in $C([-1,1])$ identifies the actual transition
law with $k_s(v,w)\varpi(dw)$. Thus its density relative to $\pi$
is uniformly bounded for $s\ge u_0$. Finally,
\[
 \int\sin^{-a_0}\theta\,\pi(d\theta)
      =c_\kappa\int_0^\pi\sin\theta\,d\theta<\infty.
\]
Markov's inequality applied to $\sin^{-a_0}\theta$ proves the last
assertion.  This is a direct proof in our normalization of the uniform
inverse-angle estimate in \cite[Lemma~2.9]{lawlerwerness2013}; see also
\cite[Lemma~2.2]{lawlerzhou2013}.
A Gegenbauer expansion for this diffusion is also used in
\cite[Proposition~4.2]{zhan2019optimal}. Here the interpolation estimate
suffices in place of explicit polynomial norm formulas.
\end{proof}

To pass from conformal radius to Euclidean clearance, write
$\delta_t=\operatorname{dist}(z,\partial D_t)$ while $t<T_z$.
This is an adapted continuous function of $t$ on the live interval.
Indeed, if $t_n\to t<T_z$, stability of the Loewner ODE gives locally
uniform convergence of the inverse maps.  Composing them with $J$ and
normalizing at $z$ gives nondegenerate conformal maps onto $D_{t_n}$
converging to the corresponding map onto $D_t$.  Thus these domains
converge in the kernel sense at $z$.  Every closed centered disk of
radius less than $\delta_t$ eventually lies in $D_{t_n}$.  Conversely,
if a fixed centered disk of radius greater than $\delta_t$ persisted
in a subsequence of the $D_{t_n}$, it would lie in the kernel $D_t$.
These two observations prove $\delta_{t_n}\to\delta_t$.

For $0<\varepsilon<\delta_0$, define the live clearance stop
\[
 S_\varepsilon(z)=\inf\{t<T_z:\delta_t\le\varepsilon\},
\]
with value infinity if this set is empty.  It is a stopping time, and
$\delta_{S_\varepsilon}=\varepsilon$ on its finite event.  This definition
also applies when $J$ has no continuous boundary extension.
When $D_t$ is the surviving domain of a continuous physical trace
$\eta$ in an initial domain $D_0$, there is the more concrete identity
\begin{equation}\label{ll:trace-clearance}
 \delta_t=\min\{\operatorname{dist}(z,\partial D_0),
                  \operatorname{dist}(z,\eta[0,t])\},\qquad t<T_z.
\end{equation}
In fact, $\partial D_t\subset\partial D_0\cup\eta[0,t]$, and both sets
on the right lie outside $D_t$; distance to the complement equals
distance to the boundary.  Thus below the initial clearance,
$S_\varepsilon$ is exactly the first live approach of the trace to
distance $\varepsilon$.  A first approach occurring only at killing
is not counted.  The finite-past domains and their similarities used
later have this continuous-trace interpretation.

\begin{lemma}[Terminal zero and approach estimates]\label{ll:approach}
Let $0<r<R_0$ and freeze the density at the first live stop with $R=r$.
Almost surely, on the event that the live radius never reaches $r$, the
density has terminal value zero,
whether its original Loewner-time lifetime is finite or infinite.
Furthermore, for $u\ge1$,
\begin{equation}
 \mathbb P\{\tau_u<\infty\}
 \le C e^{-pu}\sin^{a_0}\theta_0.
 \label{ll:radius-approach}
\end{equation}
For $0<\varepsilon\le\delta_0/(8\mathrm e)$,
\begin{equation}
 \mathbb P\{S_\varepsilon(z)<\infty\}
       \le C M_0(z)\varepsilon^p.
 \label{ll:euclidean-approach}
\end{equation}
At this stop the density has the exact first moment
\begin{equation}\label{ll:clearance-first-moment}
 \mathbb E\!\left[M_{S_\varepsilon}(z)
          \mathbf1_{\{S_\varepsilon<\infty\}}\right]=M_0(z).
\end{equation}
Under the centered law, $S_\varepsilon$ occurs no later than the finite
radius stop with $R=\varepsilon$.
All statements have their conditional restarted versions.
\end{lemma}

\begin{proof}
Suppose the radius-time lifetime is finite.  If the angle eventually
stayed in a compact subinterval of $(0,\pi)$, its SDE and
\eqref{ll:physical-clock} would extend it: the latter equations give
finite Loewner time, positive $Y$, and regular $Z$ there.  Otherwise
the angle approaches an endpoint along a sequence.  It must converge
to that endpoint.  Indeed, repeated crossings of a fixed interior
interval in a finite remaining radius time would contradict uniform
continuity of the Brownian part and boundedness of the drift on that
interval.  Thus the maximal angle solution ends at $0$ or $\pi$.
On a path that never reaches the radius $r$, its radius time is bounded
and its radius stays bounded below.  Its density therefore tends to zero,
including when its original Loewner-time lifetime is infinite.
The density stopped at the radius cutoff is bounded by $r^{-p}$;
localization followed by bounded convergence is therefore legitimate.

Under the centered law, $M_{\tau_u}/M_0$
equals $e^{pu}\sin^{a_0}\theta_u/\sin^{a_0}\theta_0$.
Untilting at this finite stop and applying
Lemma~\ref{ll:inverse-angle} gives \eqref{ll:radius-approach}.
At $S_\varepsilon<\infty$, \eqref{ll:koebe} gives
$R_{S_\varepsilon}\le4\varepsilon$.  The live radius must therefore
have crossed $8\varepsilon$ earlier.  Since
$R_0\ge\delta_0\ge8\mathrm e\varepsilon$, the corresponding radius
time is at least one, and \eqref{ll:radius-approach} proves
\eqref{ll:euclidean-approach}.

Before $S_\varepsilon$, $R\ge\delta_t>\varepsilon$
by \eqref{ll:koebe}. The stopped density is consequently bounded by
$\varepsilon^{-p}$.
Under the centered law the clearance stop occurs no later than the
finite radius level $R=\varepsilon$, because
$\operatorname{dist}(z,\partial D)\le R$.
On a path with no clearance stop, the radius stays above
$\varepsilon$, so the preceding terminal-zero assertion applies.
Bounded optional stopping and bounded convergence now give
\eqref{ll:clearance-first-moment}.  Strong Markov gives the conditional
versions with the restarted domain and its initial clearance.
\end{proof}

\begin{lemma}[Finite centered lifetime in Loewner time]\label{ll:finite-clock}
Under the consistent centered laws, $T^*(z)=\lim_{u\to\infty}\tau_u$
is finite almost surely.
\end{lemma}

\begin{proof}
The locally unique centered SDE in \eqref{ll:tilt-generator} has a global
solution $(\theta_s)_{s\ge0}$: neither endpoint is attained in finite
radius time, so the local solutions extend and agree on overlaps.
On this path reconstruct the other variables by
\[
 Y_s=Y_0\exp\left\{-\frac12\int_0^s\sin^{-2}\theta_v\,dv\right\},
 \qquad
 t_s=\int_0^s\frac{Y_v^2}{4\sin^4\theta_v}\,dv.
\]
Put $X_s=Y_s\cot\theta_s$.  Since $t_s$ is strictly increasing, a
continuous driver on every interval $[0,t_s]$ is recovered from
\[
 U_{t_s}=\operatorname{Re}\zeta+
       \int_0^s\frac{Y_v\cot\theta_v}{2\sin^2\theta_v}\,dv-X_s.
\]
This follows by integrating $\partial_t\operatorname{Re}g_t(\zeta)
=2X_t/|Z_t|^2$ and using $\operatorname{Re}g_t(\zeta)=U_t+X_t$.
Uniqueness of the centered SDE now identifies the restriction through
$s=u$ with the finite stopped law of Lemma~\ref{ll:finite-tilt}, with
$t_u=\tau_u$. This gives one probability space for the limit $T^*(z)$
and the argument below, without using a
density on an infinite-future sigma field or at the limiting time.

The formula \eqref{ll:lyapunov} at $b=a_0$ is
$\mathcal L^*V_{a_0}=pV_{a_0}$.
Thus $e^{-ps}V_{a_0}(\theta_s)$ is a nonnegative local martingale
and hence a supermartingale. Its stopped maximal inequality and
Lemma~\ref{ll:inverse-angle} imply, for integers $n\ge1$,
\[
 \mathbb P^*\left\{\sup_{n\le s\le n+1}
       V_{a_0}(\theta_s)>e^{a_0\eta n}\right\}
       \le C e^{-a_0\eta n}.
\]
For any $0<\eta<1/4$, Borel--Cantelli therefore gives
$\sin^{-4}\theta_s\le C e^{4\eta s}$ eventually.
Also \eqref{ll:physical-clock} gives $Y_s\le Y_0e^{-s/2}$.
The integral of $Y_s^2/(4\sin^4\theta_s)$ over $s\in[0,\infty)$
is finite. Initial compact intervals have finite integrals by
endpoint nonattainment. This proves the assertion solely from
finite stopped laws.
\end{proof}

\subsection{Avoidance near the real line}

\begin{lemma}[Uniform small-target estimate]\label{ll:near-real}
There are $C,c>0$ such that ordinary chordal SLE in $\mathbb H$
satisfies
\[
 \mathbb P\{\operatorname{dist}(v,\gamma)<\varepsilon\}
       \le C\varepsilon^c,
 \qquad \tfrac12\le|v|\le2,\quad
 v\in\overline{\mathbb H},\quad 0<\varepsilon<\tfrac1{16}.
\]
The same estimate holds after a simultaneous spatial rescaling of
the annulus and target radius. The trace in this statement includes
all future times and its real boundary points.
\end{lemma}

\begin{proof}
A stronger, sharp boundary exponent was proved by Alberts and
Kozdron \cite[Theorem~3.2 and Corollary~5.1]{albertskozdron2008}.
Here is a direct weaker estimate sufficient for our use.
It suffices to treat sufficiently small $\varepsilon$, since increasing
$C$ covers the remaining values below $1/16$.
Fix real $x$ with $|x|$ bounded above and away from zero, and write
$T_x$ for its real Loewner collision time. Up to $T_x$, put
$X_t=g_t(x)-\sqrt\kappa B_t$ and
$h=\int_0^tX_u^{-2}\,du$. Changing the sign of a Brownian motion
if necessary gives
\[
 \log|X|=\log|x|+(2-\kappa/2)h+\sqrt\kappa\,\widetilde B_h.
\]
For $w$ in the upper half-ball $|w-x|<4\varepsilon$, set
$D_t(w)=g_t(w)-g_t(x)$. While $|D_t(w)/X_t|\le\delta$, the ODE
gives
\[
 \frac{d}{dh}\log|D_t(w)|
      =-2\operatorname{Re}\frac{X_t}{X_t+D_t(w)}
      \le-2+C\delta.
\]
Choose $\delta>0$ so small that
$\mu=4-\kappa/2-C\delta>0$. Until the first failure of the ratio
bound, simultaneously for all these $w$,
\[
 \log\frac{|D_t(w)|}{|X_t|}
 \le\log\frac{4\varepsilon}{|x|}
             -\mu h-\sqrt\kappa\,\widetilde B_h.
\]
For sufficiently small $\varepsilon$,
$H=\log(\delta|x|/(4\varepsilon))$ is positive.  The Brownian motion
with negative drift reaches that level with probability
$\exp(-2\mu H/\kappa)$. Thus with failure
probability at most $C\varepsilon^{2\mu/\kappa}$ the ratio never
reaches $\delta$. One may first apply the argument to compact
sub-half-balls and countably many points. Analyticity and continuous
dependence extend it simultaneously to the whole open half-ball.

Put $U=B(x,4\varepsilon)\cap\mathbb H$.  On this event every point of
$U$ stays live before $T_x$, so the trace avoids $U$ there.  A
capacity-parametrized generating trace cannot stay entirely on
$\mathbb R$ during a nondegenerate time interval: adding only real
points leaves its filled half-plane hull unchanged, contrary to its
strict capacity increase.  If the trace visited
$B(x,\varepsilon)\cap\overline{\mathbb H}$ at a time $0<t\le T_x$,
continuity would place a short preceding time interval inside
$B(x,4\varepsilon)\cap\overline{\mathbb H}$.  That interval contains a
time when the trace is in $\mathbb H$, contradicting avoidance of $U$.
The initial point $0$ is outside the smaller ball.

If $T_x<\infty$, then $|X_t|\to0$, and the ratio bound shows that every
point of $U$ collides at $T_x$.  Thus $U$ is disjoint from the surviving
domain at that time.  The smaller relative half-ball is disjoint even
from its closure: each of its points has an upper-half-plane neighborhood
contained in $U$.  Future trace growth stays in that closure and cannot
visit the smaller half-ball.  If $T_x=\infty$, the preceding argument
already excludes every finite visit.  This proves the real-center estimate.

For $0\le\operatorname{Im}v\le\sqrt\varepsilon$, a visit within
$\varepsilon$ of $v$ gives a visit within $2\sqrt\varepsilon$ of
$\operatorname{Re}v$, to which the real estimate applies.
For $\operatorname{Im}v>\sqrt\varepsilon$, the radius $2\varepsilon$
is below $\delta_0/(8\mathrm e)$ once $\varepsilon$ is sufficiently
small, so \eqref{ll:euclidean-approach} applies. Here
\[
 M_0(v)=(2\operatorname{Im}v)^{-p}
             \left(\frac{\operatorname{Im}v}{|v|}\right)^{a_0}
 \le C,
\]
because $a_0-p>0$.
An actual visit to $B(v,\varepsilon)$ requires a live first hit of
$\partial B(v,2\varepsilon)$.  If the center were killed at or before
that first hit, the still untouched open ball would be entirely swallowed
and could not be entered later.  By \eqref{ll:trace-clearance}, the live
hit gives $S_{2\varepsilon}(v)<\infty$.  The approach estimate therefore
proves the result for all $v$.
\end{proof}

\subsection{A conformal localization estimate}

The next argument converts conformal area distortion into a
probability estimate. Its constants require no modulus of boundary
continuity for the conformal map. This is the uniformity needed
when a previous trace has formed narrow passages.

\begin{lemma}[Rectangle energy estimate]\label{ll:rectangle}
Let $H>0$, let $Q$ be a rectangle of fixed bounded aspect ratio, and let
$u\in W^{1,2}(Q)\cap C(\overline Q)$ satisfy $0\le u\le H$
and $u_Q\le H/2$. Write $E=\int_Q|\nabla u|^2$.
For any random set $\mathcal X$ for which every dyadic level-$k$
subrectangle of $Q$ is hit with probability at most $C_02^{-\alpha k}$,
\[
 \mathbb P\{\mathcal X\cap\{u=H\}\ne\varnothing\}
 \le C\frac E{H^2}\exp\left(-c\frac{H^2}{E}\right)
\]
whenever $E\le c_0H^2$. The constants depend only on the aspect
ratio, $C_0$, and $\alpha$. If $E=0$, the probability is zero.
\end{lemma}

\begin{proof}
We first record a two-dimensional average estimate. If $Q_k$ is
any nested chain of dyadic subrectangles with $Q_0=Q$, then
\begin{equation}
 |u_{Q_k}-u_Q|\le C\sqrt{(k+1)E}.
 \label{ll:nested-means}
\end{equation}
One proof is to reflect $u$ across the sides of $Q$ and use
concentric annuli of radii comparable to $2^{-j}$ about the center
of $Q_k$. The difference between means on consecutive annuli is
at most $C$ times the square root of the energy on their union,
by the scale-invariant Poincar\'e inequality. These unions have
bounded overlap. Cauchy--Schwarz over the $k+O(1)$ annuli gives
\eqref{ll:nested-means}. Poincar\'e on a ball of comparable radius
compares the innermost annular mean with $u_{Q_k}$ and the outermost
mean with $u_Q$. The reflected region uses a bounded number of
copies of $Q$, so its energy is at most $CE$.

Choose $N=\lfloor c_1H^2/E\rfloor$ with $c_1>0$ sufficiently small.
Decreasing $c_0$ if needed ensures $N\ge1$ and, by
\eqref{ll:nested-means}, every mean through level $N$ is at most
$3H/4$. If $u(x)=H$, choose a nested chain containing $x$.
Continuity implies that its means tend to $H$. Hence, for some
$j\ge1$, the mean changes between levels $N+j-1$ and $N+j$ by at
least $c_2H/j^2$, where $c_2$ is a sufficiently small numerical
constant. Otherwise the total change is less than $H/4$.
Poincar\'e on the parent rectangle bounds such a mean difference
by $C(\int_{Q_{N+j-1}}|\nabla u|^2)^{1/2}$.
Since level-$N+j-1$ parents have disjoint interiors, at most
$CEj^4/H^2$ of them can have this property.
They cover $\{u=H\}$ after taking the union over $j$.
The assumed hitting bound and a union bound give
\[
 C\frac E{H^2}\sum_{j\ge1}j^4
                   2^{-\alpha(N+j-1)}
 \le C\frac E{H^2}e^{-cH^2/E}.
\]
Dyadic boundary points can be assigned either adjoining chain; closed
rectangles and the stated hitting bound cover them as well.
\end{proof}

\begin{proposition}[Uniform conformal localization]\label{ll:localization}
Let $F:\mathbb H\to\mathbb C$ be univalent and admit a continuous
extension to $\overline{\mathbb H}$ with values in the sphere.
Suppose a continuous path $\eta$ in $\overline{\mathbb H}$ starts
at $0$, reaches modulus $1$, and satisfies
$F(\eta)\subset\overline B(q,r)$, where $r>0$.
There exist $b,c,C>0$ and $L_0>1$, depending only on $\kappa$, such
that for every $L\ge L_0$,
ordinary chordal SLE satisfies
\[
 \mathbb P\left\{F(\gamma[0,\sigma_{L^b}])
                      \not\subset B(q,Lr)\right\}
       \le CL^{-c},
 \qquad
 \sigma_A=\inf\{t:|\gamma(t)|=A\}.
\]
If the modulus stop is not reached, use the entire future trace.
All four constants are independent of $F,q,r$, and of the path $\eta$.
Boundary excursions are included.
\end{proposition}

\begin{proof}
Put $H=\log L$ and use strip coordinates $z=e^{v+i\phi}$,
$0<\phi<\pi$. Set
\[
 U(v,\phi)=\min\{H,\log^+(|F(e^{v+i\phi})-q|/r)\}.
\]
This function is continuous on each closed finite strip rectangle,
including at a boundary pole of $F$, where it equals $H$.
Conformal invariance of Dirichlet energy and univalence give
\begin{equation}
 \int_{\mathbb R\times(0,\pi)}|\nabla U|^2
 \le\int_{r<|w-q|<Lr}\frac{dA(w)}{|w-q|^2}
 =2\pi H.
 \label{ll:total-energy}
\end{equation}
For every $v<0$, the path assumption supplies a point on that
vertical section where $U=0$. One-dimensional Poincar\'e therefore
bounds the mean of $U$ on
$Q_j=[j,j+1]\times[0,\pi]$, $j<0$, by $C\sqrt{E_j}$, where
$E_j=\int_{Q_j}|\nabla U|^2$.

Let $A(v)=\pi^{-1}\int_0^\pi U(v,\phi)\,d\phi$.
For some $v_0\in[-1,0]$, $A(v_0)\le C\sqrt H$.
Cauchy--Schwarz applied to $A'$ gives, for $v\ge0$,
\[
 A(v)\le C\sqrt H+C\sqrt{(v+1)H}.
\]
Choose a fixed $b>0$ small enough.  There is then $H_0$, depending only
on the constants above, such that for $H\ge H_0$,
the mean on every slab with $j\le bH$ is at most $H/2$.

Under the exponential map, a dyadic level-$k$ rectangle in $Q_j$
lies in a Euclidean disk of radius $C2^{-k}e^j$ whose center has
modulus comparable to $e^j$. Lemma~\ref{ll:near-real} and Brownian
scaling therefore bound its probability of being hit by the full
SLE trace by $C2^{-\alpha k}$, uniformly in $j$.
After increasing $H_0$ if necessary, \eqref{ll:total-energy} also gives
$E_j\le2\pi H\le c_0H^2$.
Lemma~\ref{ll:rectangle} consequently bounds the probability of a
hit on $\{U=H\}$ in $Q_j$ by
$C(E_j/H^2)e^{-cH^2/E_j}$.
The infinitely many slabs to the left cause no loss: since
$E_j\le2\pi H$ and $\sum_jE_j\le2\pi H$,
\[
 \sum_{j\le bH}C\frac{E_j}{H^2}e^{-cH^2/E_j}
 \le\frac C H e^{-c'H}\le CL^{-c'}.
\]
Set $L_0=\exp(H_0)$. Before $\sigma_{L^b}$ the trace has $v\le bH$.
If its image leaves
$B(q,Lr)$, continuity gives a point with $U=H$ in one of these
slabs. The start itself maps into $\overline B(q,r)$.
This proves the result, including points on either strip edge.
\end{proof}

\subsection{Swallowing and a rule in physical disks}

\begin{lemma}[Uniform swallowing before a large excursion]
\label{ll:swallowing}
Suppose $4<\kappa<8$. For ordinary chordal SLE and every
$z\in\mathbb H$ with $|z|=1$, let $T_z$ be its swallowing time.
There are $C,c>0$ such that
\[
 \mathbb P\{\sigma_A\le T_z\}\le CA^{-c},\qquad A\ge2,
\]
uniformly over the angle of $z$.
\end{lemma}

\begin{proof}
Use the clock $h=-\log(Y/Y_0)$ and put $P=X/Y$.
Its generator is
\[
 \mathcal G=\frac\kappa4(1+P^2)\partial_P^2+2P\partial_P,
 \qquad
 \frac{dt}{dh}=\frac{Y_0^2}{2}e^{-2h}(1+P_h^2).
\]
Choose $0<\beta<1-4/\kappa$ and set
$V(P)=(1+P^2)^{\beta/2}$. Direct differentiation yields
\[
 \frac{\mathcal GV}{V}
 =\frac{2\beta P^2+(\kappa\beta/4)
                          [1+(\beta-1)P^2]}{1+P^2}.
\]
Its limit at infinity is
$\rho=2\beta+\kappa\beta(\beta-1)/4<\beta$.
Choose $\rho<\lambda<\beta$. Then
$\mathcal GV\le\lambda V+C$, and localized It\^o followed by
Gronwall gives
\begin{equation}
 \mathbb EV(P_h)\le Ce^{\lambda h}V(P_0).
 \label{ll:swallow-moment}
\end{equation}
The coefficients of the SDE for $P$ have linear growth, so its
solution exists for every finite $h$. Every such $h$ is reached by
the original time change before its lifetime: a finite terminal
$h$ would leave $P$ finite and $Y=Y_0e^{-h}>0$, while the displayed
clock gives finite Loewner time. The original Loewner equation
would then extend beyond its supposed lifetime, a contradiction.
Both drift and diffusion coefficients of $\log V(P_h)$ are bounded.
The exponential maximal inequality for its martingale part on a
unit interval therefore gives
\begin{equation}
 \mathbb E\left[\sup_{n\le h\le n+1}V(P_h)\right]
       \le Ce^{\lambda n}V(P_0).
 \label{ll:swallow-sup}
\end{equation}
One can see the conditional version directly: the martingale part
has bracket at most a constant on that interval, hence its maximum
has uniformly bounded exponential moment; the bounded drift only
changes the multiplicative constant.

Because $0<\beta/2<1$, subadditivity of $x^{\beta/2}$ and
\eqref{ll:swallow-sup} show that the terminal Loewner clock satisfies
\begin{align*}
 \mathbb E T_\infty^{\beta/2}
 &\le C Y_0^\beta\sum_{n\ge0}e^{-\beta n}
          \mathbb E\sup_{n\le h\le n+1}V(P_h)\\
 &\le C Y_0^\beta V(P_0)
             \sum_{n\ge0}e^{-(\beta-\lambda)n}
 \le C|z|^\beta=C.
\end{align*}
This finite clock ends at collision. Indeed $Y_h=Y_0e^{-h}\to0$,
and the same unit-interval estimate gives
\[
 \mathbb E\sup_{n\le h\le n+1}|X_h|^\beta
        \le C|z|^\beta e^{-(\beta-\lambda)n}.
\]
Markov and Borel--Cantelli imply $X_h\to0$.
Thus $T_\infty=T_z$, rather than a spurious endpoint of the time
change, and $\mathbb P\{T_z>A\}\le CA^{-\beta/2}$.

A hull driven through time $t$ lies below height $2\sqrt t$ and
between the minimum and maximum of its driver. The height bound
follows from $d(Y^2)/dt\ge-4$ for a point not yet swallowed;
the horizontal bound follows because a real part beyond the driver
range moves away from that range under the ODE. Consequently, for
large $A$, a trace excursion to modulus $A$ by time $A$ forces
$\sup_{t\le A}|\sqrt\kappa B_t|\ge cA$.
The Brownian reflection bound makes the probability of this event
at most $Ce^{-c'A}$. Splitting according to $T_z>A$ proves the
claim, after changing constants for bounded $A$.
\end{proof}

\begin{proposition}[The conditional physical-disk rule]
\label{ll:disk-rule}
Suppose $4<\kappa<8$. There exist $C,c>0$ and $L_0\ge2$, depending only
on $\kappa$, with the following property. Consider ordinary chordal SLE
in $\mathbb H$,
possibly after a finite stopping time with a past already drawn,
or a fixed similarity of such an evolution. Suppose
$\overline B(z,r)$ is strictly inside the current surviving domain.
Let $\tau$ be the subsequent first hit of $\partial B(z,r)$.
For $L\ge L_0$, on $\{\tau<\infty,\ z\in D_\tau\}$, set
\[
 \sigma=\inf\{t\ge\tau:|\gamma(t)-z|\ge Lr\},
\]
with value infinity if there is no such time. Then
\[
 \mathbb P\{\sigma<\infty,\ z\in D_\sigma
                    \mid\mathcal F_\tau\}\le CL^{-c}.
\]
The event in the display is an exit while the point is still live.
The constants are independent of the past, the initial disk, and the exit
time.
\end{proposition}

\begin{proof}
The first hit is a fresh interior tip. By
Lemma~\ref{geom:first-hit-access}, the radial physical segment from
this tip to $z$ has its interior in $D_\tau$ and approaches the
tip's prime end. Map $D_\tau$ to $\mathbb H$, sending tip to zero,
target to infinity, and normalize the image of $z$ to have modulus
one. Lemma~\ref{geom:extension} gives the continuous spherical
boundary extension of its inverse $F$. The radial segment supplies
exactly the path assumption of Proposition~\ref{ll:localization},
with the physical ball $\overline B(z,r)$.

Increase $L_0$ so that Proposition~\ref{ll:localization} applies and
$L^b\ge2$ for $L\ge L_0$.  Except on an event of conditional probability
$CL^{-c_1}$, the
mapped future trace stays in $B(z,Lr)$ up to the parameter-plane
modulus $L^b$. Except on an event of conditional probability
$C(L^b)^{-c_2}$, Lemma~\ref{ll:swallowing} swallows the normalized
marked point before that modulus is reached. Outside their union,
the physical point is swallowed before an exit from $B(z,Lr)$.
The domain Markov property supplies the same ordinary half-plane
law for every past. This proves the stated uniform bound.
\end{proof}
